%% ma: L12-B05-0009
%% lop: 12
%% chuong: 1
%% bai: 5
%% dang: 12.5.6
%% loai: TLN
%% muc: VDC
%% dap_an: "98,8"
%% nguon: "(NBV)-ĐỀ SỐ 7-MỖI NGÀY 1 ĐỀ THI 2025.pdf (D07, MỖI NGÀY 1 ĐỀ - Đề số 7), Phần III Câu 6"
%% kien_thuc: "Bài 5 (tối ưu thể tích bằng đạo hàm), thể tích lăng trụ, diện tích lục giác đều (Lớp 9)"
%% trang_thai: cho-duyet
%% ly_do: "Dạng toán mới đề xuất 12.5.6 (Tối ưu thể tích lăng trụ lục giác đều gấp từ tấm tôn hình tam giác đều); chờ giáo viên duyệt dạng."
%% ghi_chu: "Mức VDC. Hình dựng lại từ hình gốc: tam giác đều có viền x (hình 1), lục giác và các hình chữ nhật nét đứt (hình 2), hình khai triển (hình 3), lăng trụ (hình 4); không theo tỉ lệ tuyệt đối. Đáp án gốc 98,8 đúng (sympy: 8/81 m3 = 98,765 dm3); bảng đáp án gốc ghi nhầm 98,8 ở câu 5 và 50 ở câu 6 (hoán đổi). Điều kiện 0<x<2/3 của đề gốc không chặt: để cạnh lục giác dương cần x<√3/3≈0,577; điểm cực trị x=√3/9≈0,192 vẫn thuộc khoảng."
\begin{ex}
Cho một tấm tôn hình một tam giác đều có cạnh bằng $2$ m. Người ta thiết kế một hình lục giác đều và sáu hình chữ nhật ở phía ngoài lục giác có một cạnh bằng cạnh của lục giác, một cạnh bằng $x$ (mét) với $0<x<\dfrac23$. Sau đó người ta cắt theo nét đứt đoạn để thu được một hình hợp bởi một lục giác đều và sáu hình chữ nhật. Cuối cùng gấp các hình chữ nhật để tạo thành khối lăng trụ lục giác đều (tham khảo hình vẽ dưới đây). Thể tích của khối lăng trụ lớn nhất bằng bao nhiêu decimét khối ($\text{dm}^3$) (làm tròn kết quả đến hàng phần mười)?
\begin{center}
\begin{tikzpicture}[x=1.5cm,y=1.5cm]
  \def\r{0.5774}\def\R{1.1547}\def\xx{0.1925}\def\a{0.4444}
  % hình 1
  \begin{scope}
    \draw (-1,\r)--(1,\r)--(0,-\R)--cycle;
    \draw (-0.6667,0.3849)--(0.6667,0.3849)--(0,-0.7698)--cycle;
    \fill (0,0) circle (.8pt);
    \node[font=\scriptsize] at (0,0.481) {$x$};
  \end{scope}
  % hình 2
  \begin{scope}[xshift=3.6cm]
    \draw (-1,\r)--(1,\r)--(0,-\R)--cycle;
    \draw (\a,0)--(60:\a)--(120:\a)--(180:\a)--(240:\a)--(300:\a)--cycle;
    \foreach \k in {0,...,5}{
      \pgfmathsetmacro{\t}{60*\k}
      \pgfmathsetmacro{\ax}{\a*cos(\t)}\pgfmathsetmacro{\ay}{\a*sin(\t)}
      \pgfmathsetmacro{\bx}{\a*cos(\t+60)}\pgfmathsetmacro{\by}{\a*sin(\t+60)}
      \pgfmathsetmacro{\nx}{\xx*cos(\t+30)}\pgfmathsetmacro{\ny}{\xx*sin(\t+30)}
      \draw[dotted] (\ax,\ay)--({\ax+\nx},{\ay+\ny})--({\bx+\nx},{\by+\ny})--(\bx,\by);
    }
    \fill (0,0) circle (.8pt);
  \end{scope}
  % hình 3: hình khai triển
  \begin{scope}[yshift=-3.3cm]
    \draw (\a,0)--(60:\a)--(120:\a)--(180:\a)--(240:\a)--(300:\a)--cycle;
    \foreach \k in {0,...,5}{
      \pgfmathsetmacro{\t}{60*\k}
      \pgfmathsetmacro{\ax}{\a*cos(\t)}\pgfmathsetmacro{\ay}{\a*sin(\t)}
      \pgfmathsetmacro{\bx}{\a*cos(\t+60)}\pgfmathsetmacro{\by}{\a*sin(\t+60)}
      \pgfmathsetmacro{\nx}{\xx*cos(\t+30)}\pgfmathsetmacro{\ny}{\xx*sin(\t+30)}
      \draw (\ax,\ay)--({\ax+\nx},{\ay+\ny})--({\bx+\nx},{\by+\ny})--(\bx,\by);
    }
    \fill (0,0) circle (.8pt);
    \node[font=\scriptsize] at (0,0.48) {$x$};
  \end{scope}
  % hình 4: lăng trụ lục giác đều
  \begin{scope}[xshift=3.6cm,yshift=-3.3cm,scale=1.3]
    \def\hh{0.5}
    \foreach \k in {0,...,5}{
      \pgfmathsetmacro{\t}{60*\k}
      \pgfmathsetmacro{\px}{\a*cos(\t)}\pgfmathsetmacro{\py}{0.45*\a*sin(\t)}
      \pgfmathsetmacro{\qx}{\a*cos(\t+60)}\pgfmathsetmacro{\qy}{0.45*\a*sin(\t+60)}
      \draw ({\px},{\py+\hh})--({\qx},{\qy+\hh});
      \ifnum\k<3
        \draw[khuat] ({\px},{\py})--({\qx},{\qy});
      \else
        \draw ({\px},{\py})--({\qx},{\qy});
      \fi
    }
    \foreach \k in {0,3,4,5}{
      \pgfmathsetmacro{\t}{60*\k}
      \draw ({\a*cos(\t)},{0.45*\a*sin(\t)})--({\a*cos(\t)},{0.45*\a*sin(\t)+\hh});
    }
    \foreach \k in {1,2}{
      \pgfmathsetmacro{\t}{60*\k}
      \draw[khuat] ({\a*cos(\t)},{0.45*\a*sin(\t)})--({\a*cos(\t)},{0.45*\a*sin(\t)+\hh});
    }
  \end{scope}
\end{tikzpicture}
\end{center}
\shortans{98,8}
\loigiai{Tam giác đều cạnh $2$ có bán kính đường tròn nội tiếp $\dfrac{\sqrt3}3$. Ba cạnh của lục giác đều (cạnh $a$) nằm trên ba cạnh của tam giác đều nhỏ cách các cạnh của tấm tôn đoạn $x$, nên khoảng cách từ tâm đến cạnh của lục giác bằng $\dfrac{\sqrt3}3-x$. Do đó $\dfrac{\sqrt3}2a=\dfrac{\sqrt3}3-x\Rightarrow a=\dfrac{2-2\sqrt3x}3$.

Lăng trụ có đáy là lục giác đều cạnh $a$ và chiều cao $x$ nên
\[V(x)=\dfrac{3\sqrt3}2a^2x=\dfrac{2\sqrt3}3\,x\left(1-\sqrt3x\right)^2.\]
$V'(x)=\dfrac{2\sqrt3}3\left(1-\sqrt3x\right)\left(1-3\sqrt3x\right)=0\Leftrightarrow x=\dfrac{\sqrt3}9$ (nhận, vì $0<x<\dfrac{\sqrt3}3$; nghiệm $x=\dfrac{\sqrt3}3$ bị loại). $V'$ đổi dấu từ $+$ sang $-$ tại $x=\dfrac{\sqrt3}9$ nên $V_{\max}=V\left(\dfrac{\sqrt3}9\right)=\dfrac8{81}\ \text{m}^3\approx98{,}8\ \text{dm}^3$.}
\end{ex}
